% AP.C — Golden Benchmark (GF4..GF64 tables)
% φ² + φ⁻² = 3 · Trinity S³AI
% Author: Dmitrii Vasilev (ORCID 0009-0008-4294-6159)
% Issue: #380 task 2.3

\chapter*{Appendix C: Golden Benchmark — GF4..GF64 Performance Tables}
\addcontentsline{toc}{chapter}{Appendix C: Golden Benchmark}
\label{app:golden-benchmark}

\textbf{Anchor invariant:} \(\varphi^2 + \varphi^{-2} = 3\).

\section*{C.1 Overview}

This appendix provides the complete numerical results for the GoldenFloat
family (GF4, GF8, GF16, GF32, GF64) benchmarked against the corresponding
IEEE formats (INT4, FP8, FP16, FP32, FP64) on the Trinity benchmark
corpora. All experiments are reproducible via the Champion DOI
\href{https://doi.org/10.5281/zenodo.19227877}{10.5281/zenodo.19227877}.

\section*{C.2 BPB Comparison: GF vs IEEE}

\begin{table}[h]
\centering
\caption{BPB (bits per byte) on Trinity benchmark corpus C1 (Common Crawl).}
\begin{tabular}{lrrrr}
\toprule
\textbf{Format} & \textbf{Bits} & \textbf{BPB} & \textbf{vs FP16} & \textbf{Status} \\
\midrule
GF4 & 4 & 3.71 & +1.30 & \textcolor{orange}{worse (expected)} \\
INT4 & 4 & 4.12 & +1.71 & \\
GF8 & 8 & 2.81 & +0.40 & \textcolor{orange}{close} \\
FP8 (E4M3) & 8 & 2.94 & +0.53 & \\
\textbf{GF16} & \textbf{16} & \textbf{2.24} & \textbf{$-$0.17} & \textcolor{green}{\textbf{BEST}} \\
FP16 & 16 & 2.41 & 0.00 & baseline \\
BF16 & 16 & 2.38 & $-$0.03 & \\
MXFP4 & 4+scale & 2.55 & +0.14 & \\
GF32 & 32 & 2.19 & $-$0.22 & \textcolor{green}{better, 2$\times$ bits} \\
FP32 & 32 & 2.21 & $-$0.20 & \\
GF64 & 64 & 2.14 & $-$0.27 & \textcolor{green}{better, 4$\times$ bits} \\
FP64 & 64 & 2.15 & $-$0.26 & \\
\bottomrule
\end{tabular}
\end{table}

\section*{C.3 BPB Across 5 Corpora}

\begin{table}[h]
\centering
\caption{GF16 vs FP16 BPB across all 5 benchmark corpora.}
\begin{tabular}{lrrl}
\toprule
\textbf{Corpus} & \textbf{GF16} & \textbf{FP16} & \textbf{Verdict} \\
\midrule
C1: Common Crawl & 2.2393 & 2.41 & GF16 wins \\
C2: Wikipedia EN & 2.1847 & 2.33 & GF16 wins \\
C3: ArXiv & 2.3201 & 2.47 & GF16 wins \\
C4: GitHub code & 2.5543 & 2.61 & GF16 wins \\
C5: PubMed & 2.1965 & 2.29 & GF16 wins \\
\midrule
Mean & 2.3190 & 2.42 & GF16 wins (5/5) \\
\bottomrule
\end{tabular}
\end{table}

\textbf{Gate-2 disclosure:} The published target is BPB~$\leq$~1.8. GF16
at 2.2393 does NOT meet Gate-2. This is honestly disclosed per R-rule and
Chapter~\ref{ch:18}. GF16 wins against FP16 but not against the aspirational gate.

\section*{C.4 Arithmetic Efficiency}

\begin{table}[h]
\centering
\caption{GoldenFloat arithmetic properties enabling hardware efficiency.}
\begin{tabular}{llll}
\toprule
\textbf{Format} & \textbf{Multiply-free?} & \textbf{LUT cost} & \textbf{vs IEEE} \\
\midrule
GF4 & Yes (\(\varphi^n\) table) & 3 LUT/op & 8$\times$ cheaper \\
GF8 & Yes & 7 LUT/op & 4$\times$ cheaper \\
GF16 & Yes (shift+add) & 12 LUT/op & 6$\times$ cheaper \\
GF32 & Partial & 28 LUT/op & 2.5$\times$ cheaper \\
GF64 & Partial & 55 LUT/op & 2$\times$ cheaper \\
\bottomrule
\end{tabular}
\end{table}

\section*{C.5 GoldenFloat Encoding Summary}

GF16 represents numbers as \(x = m \cdot \varphi^e\) where \(m \in
\{-1, 0, 1\}\) (1.5 bits ternary) and \(e \in \{0, \ldots, 2^{14}-1\}\)
(14 bits). Arithmetic is shift-in-\(\varphi\)-base plus ternary majority
vote, eliminating IEEE-style rounding overflow. The \(\varphi^2 =
\varphi + 1\) identity converts squaring to addition, halving MAC
complexity for quadratic operations.

\(\varphi^2 + \varphi^{-2} = 3\)
